% Copyright 2026 Nobin Sijo (NobinSijo7T).
% SPDX-License-Identifier: Apache-2.0
\documentclass[11pt,a4paper]{article}
\usepackage[margin=0.68in,top=0.56in,bottom=0.56in]{geometry}
\usepackage{amsmath,amssymb}
\usepackage{times}
\usepackage{enumitem}
\usepackage{fancyhdr}

\pagestyle{fancy}
\fancyhf{}
\fancyhead[L]{\small Group Project $\cdot$ PrismSpace Architecture}
\fancyhead[R]{\small\thepage}
\renewcommand{\headrulewidth}{0.4pt}
\setlength{\headheight}{14.5pt}

\setcounter{page}{6}

\setlength{\parindent}{0pt}
\setlength{\parskip}{0.18em}
\setlength{\abovedisplayskip}{3pt}
\setlength{\belowdisplayskip}{3pt}
\setlength{\abovedisplayshortskip}{2pt}
\setlength{\belowdisplayshortskip}{2pt}

\begin{document}

% =========================================================================
% PAGE 1 (Page 6): Overview, Inputs, Outputs, and Phase I (Offline Graph & Routing)
% =========================================================================

\subsection*{4.5\quad Algorithm of PrismSpace Multi-Agent System}

The PrismSpace algorithm orchestrates specialized autonomous AI agents into an adaptive collaborative swarm: parsing complex user objectives, routing sub-tasks to cost- and latency-optimal LLM providers, dispatching specialized worker agents equipped with Model Context Protocol (MCP) tools, and enforcing rigorous preference alignment and safety gating. The algorithm operates across three cohesive phases:
\begin{itemize}[leftmargin=1.5em, itemsep=0.04em, topsep=0.03em]
    \item \textbf{Phase I: Offline Swarm Graph Representation \& Routing Optimization} (Steps 1--3: Swarm graph topology, Graph Attention Network embedding, multi-task routing and dispatch loss).
    \item \textbf{Phase II: Single-Stage Preference Alignment} (Step 4: Unified SFT and Odds-Ratio preference tuning via ORPO without external reward models).
    \item \textbf{Phase III: Online Runtime DAG Orchestration \& Swarm Execution} (Steps 5--7: Task subsection division, DAG construction, topological scheduling, ReAct tool execution, safety gating, and telemetry).
\end{itemize}

\vspace{0.10em}
\textbf{System Inputs \& Outputs:}
\begin{itemize}[leftmargin=1.5em, itemsep=0.03em, topsep=0.02em]
    \item \textbf{Inputs:} User objective prompt $\mathcal{T}_{\text{user}}$, curated training corpora $\mathcal{D}$ (intent, routing, and safety traces across RouterBench, LMSYS, WildGuard, CMU Agent benchmarks), preference pairs $\mathcal{D}_{\text{pref}} = \{(x, y_w, y_l)\}$, and global agent/tool registry $\mathcal{G}_{\text{global}}$.
    \item \textbf{Outputs:} Dispatched agent team $\hat{\mathbf{y}}_{\text{agents}}$, optimal LLM provider assignment $\hat{y}_{\text{provider}}$, execution DAG $\mathcal{G}_{\text{active}}$, verified streaming deliverables, and persistent session state in browser IndexedDB (\texttt{PrismDB}).
\end{itemize}

\vspace{0.15em}
\textbf{Phase I: Offline Swarm Graph Representation \& Intelligent Routing}

\textbf{Step 1: Swarm Graph Formulation ($\mathcal{G} = (V, E)$)}
\begin{enumerate}[leftmargin=1.8em, itemsep=0.06em, topsep=0.02em]
    \item \textit{Node Definition ($V$):} Represent each specialized agent persona (Supervisor, Research, Coding, Validator, Planner) and each MCP tool capability (Browser, Filesystem, GitHub, Sandbox CLI) as a distinct vertex $v_i \in V$. Each node is initialized with a normalized capability feature vector $\mathbf{h}_i \in \mathbb{R}^d$.
    \item \textit{Collaborative Edge Topology ($E$):} Establish undirected communication edges $e_{ij} \in E$ between complementary agent-tool and agent-agent pairs whose functional correlation exceeds an empirical interaction threshold $\tau$, permitting direct message passing.
\end{enumerate}

\vspace{0.08em}
\textbf{Step 2: Attention-Based Message Passing (Graph Attention Network)}
To dynamically model swarm affinity and contextual synergy, node embeddings are updated using multi-head graph attention:
\begin{equation}
\alpha_{ij}^k = \frac{\exp\Big(\mathrm{LeakyReLU}\big((\mathbf{a}^k)^T [\mathbf{W}^k \mathbf{h}_i \,\|\, \mathbf{W}^k \mathbf{h}_j]\big)\Big)}{\sum_{m \in \mathcal{N}(i)} \exp\Big(\mathrm{LeakyReLU}\big((\mathbf{a}^k)^T [\mathbf{W}^k \mathbf{h}_i \,\|\, \mathbf{W}^k \mathbf{h}_m]\big)\Big)}, \qquad \mathbf{h}_i' = \Vert_{k=1}^K \sigma\left(\sum_{j \in \mathcal{N}(i)} \alpha_{ij}^k \mathbf{W}^k \mathbf{h}_j\right)
\end{equation}
\textit{Notation:} $\mathbf{h}_i \in \mathbb{R}^d$ is the node feature vector; $\mathcal{N}(i)$ is the set of collaborative neighbors of node $v_i$; $\mathbf{W}^k \in \mathbb{R}^{d' \times d}$ is the head-specific linear transformation matrix; $\mathbf{a}^k \in \mathbb{R}^{2d'}$ is the parameterized attention mechanism vector for head $k \in \{1,\dots,K\}$; $\|$ denotes vector concatenation; $\alpha_{ij}^k$ is the normalized attention coefficient reflecting how critical node $v_j$ is to node $v_i$; and $\mathbf{h}_i' \in \mathbb{R}^{K \cdot d'}$ is the updated context-aware node embedding.

\vspace{0.08em}
\textbf{Step 3: Multi-Task Joint Routing and Agent Dispatch Objective}
The pooled swarm representation $\mathbf{h}'$ and individual agent states $\mathbf{h}_i'$ are projected through dual prediction heads:
\begin{enumerate}[leftmargin=1.8em, itemsep=0.06em, topsep=0.02em]
    \item \textit{Provider Routing Head:} Computes posterior probability distribution $\hat{\mathbf{y}}_{\text{route}} = \mathrm{Softmax}(\mathbf{W}_{\text{route}} \mathbf{h}' + \mathbf{b}_{\text{route}})$ over available LLM providers (Groq for ultra-low latency, OpenAI GPT-4o for deep reasoning, Claude 3.5 Sonnet for code).
    \item \textit{Agent Dispatch Head:} Computes independent activation probabilities $\hat{y}_v = \mathrm{Sigmoid}(\mathbf{w}_{\text{agent}}^T \mathbf{h}_v' + b_v)$ for each candidate agent persona $v \in V$.
    \item \textit{Joint Optimization Loss:} Parameters are trained end-to-end via AdamW using the multi-task objective:
    \begin{equation}
    \mathcal{L}_{\text{class}} = \mathcal{L}_{\text{CE}}(\mathbf{y}_{\text{route}}, \hat{\mathbf{y}}_{\text{route}}) + \lambda \sum_{v \in V} \mathcal{L}_{\text{BCE}}(y_v, \hat{y}_v)
    \end{equation}
    where $\mathcal{L}_{\text{CE}}$ is categorical cross-entropy penalizing incorrect LLM provider routing against ground truth $\mathbf{y}_{\text{route}}$, $\mathcal{L}_{\text{BCE}}$ is binary cross-entropy penalizing over- or under-provisioning of agents against optimal ground truth $y_v \in \{0, 1\}$, and $\lambda$ is a balancing hyperparameter.
\end{enumerate}

\clearpage

% =========================================================================
% PAGE 2 (Page 7): Phase II (ORPO Alignment) & Phase III (Task Subsections & DAG Graph)
% =========================================================================

\textbf{Phase II: Single-Stage Preference Alignment (ORPO)}

\textbf{Step 4: Odds Ratio Preference Optimization}
Rather than relying on unstable multi-stage Reinforcement Learning from Human Feedback (RLHF) requiring a separate reward model and PPO policy updates, PrismSpace tunes the agent models in a single training stage via Odds Ratio Preference Optimization (ORPO) paired with Low-Rank Adaptation (LoRA, $\Delta \mathbf{W} = \frac{\alpha}{r} \mathbf{B}\mathbf{A}$):
\begin{equation}
\mathcal{L}_{\mathrm{ORPO}} = \mathcal{L}_{\mathrm{SFT}} + \beta \cdot \mathcal{L}_{\mathrm{OR}}, \qquad \mathcal{L}_{\mathrm{OR}} = -\log \sigma \left(\log \left[\frac{\mathrm{odds}(y_w \mid x)}{\mathrm{odds}(y_l \mid x)}\right]\right)
\end{equation}
\textit{Notation:} $x$ is the input instruction prompt; $y_w$ is the winning/favored response; $y_l$ is the losing/disfavored response (e.g., hallucinated facts, ungrounded actions, safety policy breaches); $\mathcal{L}_{\mathrm{SFT}} = -\frac{1}{|y_w|}\sum_{t=1}^{|y_w|} \log P(y_{w,t} \mid x, y_{w,<t})$ is the standard supervised causal language modeling loss maintaining fluency; $\mathrm{odds}(y \mid x) = \frac{P(y \mid x)}{1 - P(y \mid x)}$ calculates the length-normalized sequence odds; $\sigma(\cdot)$ is the sigmoid function; and $\beta$ is a weighting coefficient penalizing generation paths that favor rejected completions.

\vspace{0.18em}
\textbf{Phase III: Task Subsection Division \& Directed Graph Construction}

A central mechanism of the PrismSpace algorithm is decomposing an unstructured, high-level user objective $\mathcal{T}_{\text{user}}$ into an executable, dependency-aware workflow. This decomposition and graph synthesis operates as follows:

\vspace{0.06em}
\textbf{1. What Vertices ($V_{\text{task}}$) Are Made Of:}
\begin{itemize}[leftmargin=1.5em, itemsep=0.06em, topsep=0.02em]
    \item When prompt $\mathcal{T}_{\text{user}}$ is ingested, the Supervisor/Planner agent decomposes the prompt into $N$ discrete, atomic task subsections: $V_{\text{task}} = \{v_1, v_2, \dots, v_N\}$.
    \item Each vertex $v_i \in V_{\text{task}}$ represents an independent operational sub-task constructed from five concrete elements:
    \begin{equation*}
    v_i = \big(\text{TaskID}_i,\, \text{Goal}_i,\, \mathcal{I}_i,\, \mathcal{O}_i,\, \text{Role}_i,\, \text{Provider}_i,\, \text{Tools}_i\big)
    \end{equation*}
    where $\text{Goal}_i$ is the target instruction; $\mathcal{I}_i$ is the required input schema; $\mathcal{O}_i$ is the deliverable output artifact (e.g., CSV data, research summary, code patch); $\text{Role}_i \in \{\text{Planner}, \text{Researcher}, \text{Coder}, \text{Assistant}, \text{Reviewer}\}$ is the assigned agent persona; $\text{Provider}_i$ is the routed LLM backend; and $\text{Tools}_i$ is the set of granted MCP tools.
    \item This isolates concerns: each agent receives only its necessary inputs, eliminating context bloat and token waste.
\end{itemize}

\vspace{0.06em}
\textbf{2. What Each Edge ($E_{\text{dep}}$) Does \& How Vertices and Edges Connect:}
\begin{itemize}[leftmargin=1.5em, itemsep=0.06em, topsep=0.02em]
    \item \textit{How They Connect:} An edge is a directed link $e = \langle v_i, v_j \rangle \in E_{\text{dep}}$ pointing from producer vertex $v_i$ to consumer vertex $v_j$, linking the output interface of $v_i$ directly to the input interface of $v_j$.
    \item \textit{What Each Edge Does:} Each edge performs two operational functions:
    \begin{enumerate}[leftmargin=1.4em, itemsep=0.02em, topsep=0.01em]
        \item \textbf{Data Pipeline:} Automatically transmits the completed output artifact $\mathcal{O}_i$ from $v_i$ into the input buffer $\mathcal{I}_j$ of $v_j$.
        \item \textbf{Execution Gate:} Blocks task $v_j$ from starting until task $v_i$ completes its execution and delivers $\mathcal{O}_i$.
    \end{enumerate}
\end{itemize}

\vspace{0.06em}
\textbf{3. Basis for Connecting Adjacent Nodes (Data Dependency Matching):}
\begin{itemize}[leftmargin=1.5em, itemsep=0.06em, topsep=0.02em]
    \item Adjacent nodes are connected \textbf{strictly on the basis of Data Dependency}: two nodes $v_i$ and $v_j$ are linked with directed edge $\langle v_i, v_j \rangle$ \textbf{if and only if} task $v_j$ requires the output artifact generated by task $v_i$:
    \begin{equation*}
    \langle v_i, v_j \rangle \in E_{\text{dep}} \iff \mathcal{O}_i \cap \mathcal{I}_j \neq \emptyset
    \end{equation*}
    \item \textit{Concrete Connection Example:} For a request like \textit{``Fetch quarterly financial metrics, scrape market news, synthesize executive insights, and format an email draft''}:
    \begin{itemize}[leftmargin=1.4em, itemsep=0.02em, topsep=0.01em]
        \item Nodes $v_1$ (Fetch CSV) and $v_2$ (Search News) have no input dependencies $\to$ they have no incoming edges.
        \item Node $v_3$ (Synthesize Analysis) requires inputs from both $v_1$ and $v_2$ ($\mathcal{I}_3 = \{\mathcal{O}_1, \mathcal{O}_2\}$) $\to$ edges $\langle v_1, v_3 \rangle$ and $\langle v_2, v_3 \rangle$ connect them.
        \item Node $v_4$ (Draft Email) requires $\mathcal{O}_3$ ($\mathcal{I}_4 = \{\mathcal{O}_3\}$) $\to$ edge $\langle v_3, v_4 \rangle$ connects them.
    \end{itemize}
    \item \textit{Parallel Concurrency for Independent Nodes:} Tasks with no connecting dependency edges ($\text{in-degree} = 0$, such as $v_1$ and $v_2$) execute simultaneously in parallel (concurrency limit $K \le 3$). When $v_1$ and $v_2$ complete, their outgoing edges release $v_3$ for execution.
\end{itemize}

\clearpage

% =========================================================================
% PAGE 3 (Page 8): Phase IV (Online Execution & Verification) and Summary
% =========================================================================

\textbf{Phase IV: Online Swarm Execution, Tool Interaction \& Verification}

\textbf{Step 5: Autonomous ReAct Swarm Execution}
For each scheduled task subsection $v_i \in \pi$:
\begin{enumerate}[leftmargin=1.8em, itemsep=0.06em, topsep=0.02em]
    \item Dispatched agent $v_i$ instantiates an autonomous Reasoning-Action (ReAct) cycle:
    \begin{equation*}
    \text{Observation}_t \longrightarrow \text{Thought}_t \longrightarrow \text{Action}_t(\text{Tool}, \text{Params}) \longrightarrow \text{Observation}_{t+1}
    \end{equation*}
    \item The agent issues JSON-RPC calls across Model Context Protocol (MCP) server endpoints (e.g., executing Python scripts in a sandboxed CLI, writing files via Filesystem MCP, querying live APIs via Browser MCP).
    \item Context is dynamically updated: $C_{t+1} = \mathrm{Compress}(C_t \cup \Delta C_t \setminus D_t)$, evicting stale tool observations $D_t$ to prevent token window overflow.
\end{enumerate}

\vspace{0.08em}
\textbf{Step 6: Multi-Criteria Safety Gating and Anomaly Isolation}
Before any generated deliverable is released to downstream tasks or the user interface, it must pass an automated multi-criteria safety evaluation barrier:
\begin{equation}
S = \sum_{m=1}^M w_m \, q_m \ge \tau_{\text{safety}}
\end{equation}
where $q_m \in [0, 1]$ are individual verification scores across: (a) WildGuardMix adversarial prompt/harm checks, (b) IsolationForest tool parameter anomaly scores, (c) schema completeness against target specification $\mathcal{O}_i$, and (d) ORPO alignment safety margins; $w_m$ are normalized criterion weights ($\sum w_m = 1$); and $\tau_{\text{safety}}$ is the strict release threshold. Deliverables violating safety gating are redirected back to the Planner for corrective re-planning.

\vspace{0.08em}
\textbf{Step 7: Deliverable Rendering and Reactive IndexedDB Persistence}
\begin{enumerate}[leftmargin=1.8em, itemsep=0.06em, topsep=0.02em]
    \item Verified partial results, reasoning traces, and artifacts are rendered live in the Next.js frontend via reactive local database subscriptions (\texttt{useLiveQuery()}).
    \item The complete execution trace, including DAG state, agent conversation turns, tool call arguments, and final artifacts, is committed atomically to the browser's local IndexedDB instance via Dexie.js (\texttt{PrismDB}), guaranteeing zero data loss, offline accessibility, and total user privacy.
\end{enumerate}

\vspace{0.12em}
\textbf{Model Evaluation and Telemetry Diagnostics}
\begin{itemize}[leftmargin=1.5em, itemsep=0.06em, topsep=0.02em]
    \item \textit{Routing and Safety Generalization:} Evaluated across holdout test splits ($\mathcal{D}_{\text{test}}$) using Macro-F1 (routing accuracy), ROC-AUC / PR-AUC (safety discrimination), and Expected Calibration Error ($\mathrm{ECE} < 0.05$).
    \item \textit{Swarm Telemetry HUD:} Real-time operational tracking of task execution latency $\hat{\ell}$ (ms), cumulative token costs $\hat{c}$ (USD), agent utilization rates, and cache hit ratios.
\end{itemize}

\vspace{0.14em}
\textbf{Algorithm Summary \& Mathematical Guarantee}
\begin{itemize}[leftmargin=1.5em, itemsep=0.06em, topsep=0.02em]
    \item \textbf{Deterministic Execution Order:} Decomposing user goals into task subsections $V_{\text{task}}$ connected by data dependency edges $E_{\text{dep}}$ guarantees that each task receives its required input artifacts before starting.
    \item \textbf{Optimal Resource Allocation:} GAT swarm attention and multi-task loss $\mathcal{L}_{\text{class}}$ ensure that only strictly necessary agents are instantiated and routed to cost- and latency-optimal LLM providers.
    \item \textbf{Single-Stage Alignment Stability:} ORPO loss $\mathcal{L}_{\mathrm{ORPO}}$ enforces safety and human preference directly during parameter tuning without the instability, latency, or compute overhead of external reward models.
\end{itemize}

\end{document}
